\chapter{Phone and Technical Screens}\label{ch:iv:phone-and-technical-screens}

Twenty minutes into a video screen the candidate announces an expected value of 3.5 for a game whose payoff
is never above three. The interviewer says nothing. The candidate who then says ``that cannot be right: the
payoff is capped at three'' and finds the error in a minute has had a better screen than one who never made
it, and the notes will say so. A screen is short, it is decided on a handful of observations, and two of the
most valuable ones are how the candidate thinks when speaking and what the candidate does after a mistake.

\section{What a screen is for}

\begin{definition}[Technical screen]\label{def:iv:phone-and-technical-screens:screen}
A \emph{technical screen}\index{technical screen} is a short \omterm{def:iv:what-each-interview-tests-by-role:loop}{technical interview}, by telephone or video,
that decides whether a candidate goes on to the final round. It has a pass-or-fail outcome and a short
rubric, usually two to four questions of increasing difficulty.
\end{definition}

A screen asks less than a final round and forgives less. Its interviewer has thirty to forty-five minutes and
one decision to make, so the rubric records a few things: did the candidate reach answers, did the candidate
explain them, did the candidate check them, and would a colleague enjoy an afternoon of this. It does not
record how much the candidate knows beyond the questions asked. Video removes most of the non-verbal channel,
so the spoken account of the reasoning is almost all the interviewer has.

\section{Thinking aloud}

\begin{definition}[Think-aloud protocol]\label{def:iv:phone-and-technical-screens:thinkaloud}
The \emph{think-aloud protocol}\index{think-aloud protocol} is the practice of saying one's reasoning while
solving a problem, as it happens, without explaining it to the listener after the fact. The term comes from
the psychology of problem solving, where verbal reports of this kind are recorded as data about the solver's
thinking (Ericsson and Simon).
\end{definition}

The interviewer cannot score what the interviewer cannot hear. A silent minute followed by a correct number
scores less than a minute in which the candidate says what the problem is, what method applies, and what the
first step gives; and a wrong final number after a correct spoken method is often scored as a slip, not a
failure. Thinking aloud is a skill with a structure.

\begin{method}[Thinking aloud in a screen]\label{met:iv:phone-and-technical-screens:aloud}
\begin{enumerate}
\item Restate the question in one sentence, with the assumptions you will make (``fair die, independent
rolls'').
\item Name the method before using it (``first-step analysis on two states'', ``linearity with
indicators'').
\item Compute in steps small enough to say, writing the intermediate numbers where the interviewer can see
them.
\item Check (the next section) and say the check.
\item Conclude with the number and one sentence of interpretation, then stop talking.
\item Silence is allowed when announced: ``give me twenty seconds to set this up''.
\end{enumerate}
\end{method}

\section{Sanity checks}

\begin{definition}[Sanity check]\label{def:iv:phone-and-technical-screens:sanity}
A \emph{sanity check}\index{sanity check} is a fast test of an answer against a property it must have: a
bound, a unit, a sign, a symmetry, a limiting case, a small case that can be computed by hand, or an order of
magnitude from a second route. It cannot prove an answer right; it can prove an answer wrong in seconds.
\end{definition}

\begin{method}[The sanity checks, in the order to try them]\label{met:iv:phone-and-technical-screens:checks}
\begin{enumerate}
\item \emph{Range}: a probability in $[0,1]$; an expectation between the smallest and largest outcomes; a
variance non-negative; a correlation in $[-1,1]$.
\item \emph{Units and sign}: money, time and counts in the right units; a price that rises with the thing
that should raise it.
\item \emph{Limits}: set a parameter to 0, 1 or infinity and see whether the formula gives the obvious answer.
\item \emph{Small case}: $n = 1$ or $n = 2$ by hand.
\item \emph{Symmetry}: swapping two roles that should not matter must not change the answer.
\item \emph{Second route}: an order of magnitude by a different argument.
\end{enumerate}
\end{method}

\begin{example}[A range check]\label{ex:iv:phone-and-technical-screens:six}
``What is the chance of at least one six in six rolls?'' The quick answer $6 \times \tfrac16 = 1$ fails the
range check at once: the event is not certain, since all six rolls can miss. The union bound overcounts the
overlaps; the complement gives $1 - (5/6)^6 \approx 0.665$.
\end{example}

\begin{example}[A limit check]\label{ex:iv:phone-and-technical-screens:limit}
A candidate proposes $\sigma\sqrt{w_1^2 + w_2^2}$ for the volatility of a two-asset portfolio with weights
$w_1, w_2$ and equal volatilities $\sigma$. Setting the correlation to 1 must return $\sigma(w_1 + w_2)$, since the
assets are then the same; the formula does not depend on the correlation at all, so it is the answer for
correlation 0 only.
\end{example}

\section{Recovering from an error}

Errors in a screen are normal; what is scored is what happens next. Say that the check failed, say which
step you suspect, redo that step aloud, and state the corrected answer with the check that now passes. Do not
restart from the beginning in silence, and do not argue that the wrong answer was nearly right.

Asking for a hint is also allowed. It costs a little on the rubric and saves the screen: a candidate stuck for
ten minutes gives the interviewer nothing to write down. The good form is specific: ``I have set up the two
states and I am not sure whether the absorbing state counts as a step; is that the right way to think about
it?'' A vague ``can you help me?'' costs more and teaches less.

Coding screens add one more check: run the code on a small case, by hand if there is no interpreter. The
function in \cref{lst:iv:phone-and-technical-screens:windows} was submitted in a screen; the small case finds
its error (\cref{iq:iv:phone-and-technical-screens:3}).

\omcode{code/interviews/05-phone-and-technical-screens/python/iv_screen_windows.py}{4}{9}{A moving average as
submitted in a screen. Run it by hand on four prices with $k = 2$.}\label{lst:iv:phone-and-technical-screens:windows}

The last five minutes of a screen usually belong to the candidate. One or two questions that show the
candidate has thought about the role (what a first year works on, how work is reviewed, what distinguishes
people who do well) are worth more than questions whose answers are on the firm's website.

\section{Worked answers}

\begin{example}[A screen question answered aloud]\label{ex:iv:phone-and-technical-screens:product}
\emph{``A fill ratio $X$ and a price improvement factor $Y$ are independent and uniform on $[0,1]$. What is the chance
that their product exceeds one half?''}

\emph{Aloud.} ``I'll draw the unit square, $X$ across and $Y$ up. The product exceeds a half above the hyperbola $y =
1/(2x)$, which only enters the square for $x$ above a half. So the area is the integral from a half to one of $1 -
1/(2x)$, which is $\tfrac12 - \tfrac12\ln 2$, about $0.5 - 0.347 = 0.153$. Checks: for a general level $c$ the same
steps give $1 - c + c\ln c$; at $c = 0$ that is 1 and at $c = 1$ it is 0, both right; and the answer must be below
$1/4$, because a product above a half needs both numbers above a half. So 0.15 is plausible.'' Three things score beyond the number: the picture chosen out loud, the integral set up
before it is computed, and two limit checks in ten seconds.
\end{example}

\begin{example}[Recovering from a wrong first answer]\label{ex:iv:phone-and-technical-screens:recover}
\emph{``How many rolls of a die, on average, until you see a six immediately followed by a roll that is not a six?''}

\emph{Aloud.} ``Two consecutive specified rolls, so 36?'' The candidate stops: ``That's the answer for two sixes in a
row, where a failure sends me back to the start. Here it doesn't: if a six is followed by another six, I've still just
rolled a six. Let me redo it. Waiting for the first six takes 6 rolls on average. From `just rolled a six', each roll ends
the wait with probability 5/6 and otherwise leaves me in the same state, so 6/5 more rolls. Total $6 + 6/5 = 7.2$.
Check: it has to be at least 6, since I need a six first, and much less than 36, since the second roll is easy. 7.2
passes both.'' The interviewer writes down the recovery, not the false start: the failed check was named, the faulty step
identified, and the new answer checked.
\end{example}

\begin{example}[A range check on units]\label{ex:iv:phone-and-technical-screens:units}
\emph{``A stock's daily returns have a standard deviation of 1.5\%. What is its annualised volatility?''} A candidate who
multiplies by 252 says 378\% and should hear the alarm at once: a large-cap stock with a volatility of nearly four
hundred per cent a year is outside every range they know. Variances add over independent days, so the standard deviation
grows like the square root of time: $1.5\% \times \sqrt{252} \approx 1.5\% \times 15.9 \approx 23.8\%$, an ordinary
equity volatility.
\end{example}

\section{Question bank}

\begin{interviewq}[$\star$ \iqroles{trader} \iqfirm{market maker}]\label{iq:iv:phone-and-technical-screens:1}
A candidate says that the expected value of the larger of two fair dice is 3.5. Without computing, why is that
wrong? Then compute it.
\end{interviewq}

\begin{interviewq}[$\star$ \iqroles{researcher, risk} \iqfirm{systematic fund}]\label{iq:iv:phone-and-technical-screens:2}
Ten independent alerts each fire on a given day with probability 0.15. A colleague says the chance that at
least one fires is 1.5. Say what is wrong in one sentence and give the right number.
\end{interviewq}

\begin{interviewq}[$\star$ \iqroles{developer} \iqfirm{proprietary firm}]\label{iq:iv:phone-and-technical-screens:3}
Run \cref{lst:iv:phone-and-technical-screens:windows} by hand on the prices 10, 11, 12, 13 with $k = 2$. What
does it return, what should it return, and how would you fix it so that it also runs in $O(n)$ time?
\end{interviewq}

\begin{interviewq}[$\star\star$ \iqroles{trader, risk} \iqfirm{bank}]\label{iq:iv:phone-and-technical-screens:4}
Two assets each have a volatility of 20\% and their correlation is 0.5. A candidate estimates the volatility
of an equally weighted portfolio at 25\%. Which check rejects this at once? What is the right figure?
\end{interviewq}

\begin{interviewq}[$\star\star$ \iqroles{researcher, developer} \iqfirm{any}]\label{iq:iv:phone-and-technical-screens:5}
You conjecture in a screen that a $2 \times n$ board can be tiled by $2 \times 1$ dominoes in $2^{n-1}$ ways.
Test it on small cases before the interviewer does. What is the right count?
\end{interviewq}

\begin{interviewq}[$\star\star$ \iqroles{bank, researcher} \iqfirm{bank}]\label{iq:iv:phone-and-technical-screens:6}
Twelve minutes into a derivation you are stuck, and there are fifteen minutes left. What exactly do you say,
and what should you avoid?
\end{interviewq}

\begin{interviewq}[$\star\star\star$ \iqroles{trader} \iqfirm{market maker}]\label{iq:iv:phone-and-technical-screens:7}
You roll a die and may roll once more, keeping the second result if you do. A candidate argues: ``I re-roll
half the time; when I keep the first roll it averages 5; so the game is worth 4.5.'' Find the error, give the
value, and give the value when up to three rolls are allowed.
\end{interviewq}

\begin{interviewq}[$\star\star\star$ \iqroles{researcher, developer} \iqfirm{systematic fund}]\label{iq:iv:phone-and-technical-screens:8}
You state in a screen that the expected number of rolls of a fair die until every face has appeared is about
14.7. The interviewer asks how you would check that in thirty seconds without a computer. What do you say?
\end{interviewq}

\begin{omsources}
\item K.~A.~Ericsson and H.~A.~Simon, \emph{Protocol Analysis: Verbal Reports as Data}, MIT Press, 1984;
revised edition 1993.
\item G.~P\'olya, \emph{How to Solve It}, Princeton University Press, 1945, on looking back at a solution.
\end{omsources}
